\begin{abstract}
LLM agents are being built to investigate security alerts by walking provenance evidence through tools:
which process started which, what it wrote, where it connected. In this adversarial setting an agent must be
right for the right reasons, at bounded cost, on telemetry an attacker partly controls. We present an
investigation agent that routes every tool call through local safeguards (response cache, schema validation, a
grounding check on submitted identifiers, and call budgets), and \bench, a benchmark of 21 provenance-graph tasks
built from public Windows Sysmon and Sysmon-for-Linux attack recordings: 15 ATT\&CK technique executions and 6
benign look-alikes from the same hosts, including twins that share an alert with an attack.%
\ifnum\ResStudyAvailable=1
\ Over \ResStudyRepeats{} repeats with Claude Opus 5.5 the agent reached reward \ResFullReward{}, verdict
accuracy \ResFullAccuracy{} and no false positives, at \ResFullCalls{} tool calls per task. A paired ablation
shows that the investigation method in the prompt drives efficiency (\ResNoMethodCalls{} calls per task
without it, more on \ResNoMethodMoreCalls{} of \ResNoMethodCallsDiffer{} tasks), while the safeguards had no
measurable effect: per-call traces show they fired on fewer than one call in ten tasks. The traces instead
expose what aggregate scores hide: identifiers of the wrong type that a substring
grounding check accepts, an efficiency instruction that withheld the identifier the response needed, and label
cues that public attack recordings leak into the alert, which \TrAllSourceCuePct\% of submissions cite. We
release \emph{blind} and \emph{mimicry} variants of \bench that remove these cues and rename attacker-chosen
artefacts without changing behaviour, so that shortcut reliance and lexical evasion become measurable.
\else
\ We define an ablation protocol that isolates each safeguard's contribution.
\fi
Agent, benchmark, traces of every tool call and the scripts that regenerate every table are released.
\end{abstract}

\begin{IEEEkeywords}
Security operations, provenance graphs, LLM agents, tool use, robustness, shortcut learning, mimicry, benchmark, ablation study.
\end{IEEEkeywords}
